\section{The all-class commutator recursion}\label{app:commutator}

This appendix develops the candidate proof of Proposition~\ref{prop:n8}.
The untruncated rational free Lie algebra is $L=L_\Z\otimes\Q$,
graded by weight. Hall group coordinates remain integral throughout;
rational changes below prove structural separation and do not replace
the integral solution lattice. The elementary cases $g=1$, rank at
most one and class at most one are handled directly.

\subsection{Complete leading branches}
The exact moves $(x,y)\mapsto(yx,y)$ and $(x,y)\mapsto(x,xy)$
preserve $[x,y]$. The rotation
$(x,y)\mapsto(y^{-1},y^{-1}xy)$ does so as well. On a common leading
layer these realize $\SL_2(\Z)$. If the two leading terms are
dependent, Euclidean reduction kills one of them and increases its
weight. Nonzero commuting elements of an untruncated free Lie algebra
are proportional: otherwise they span a two-dimensional abelian
subalgebra, contrary to Shirshov--Witt. The normalized nonidentity
target therefore has leading branches
\[
 [C,D]=W,\quad C\in(L_\Z)_p,\quad D\in(L_\Z)_q,
 \quad p\le q,\quad p+q=d.
\]

For $p<q$ the bracket has no zero divisors, also over an algebraic
closure. Its projective linear projection on
$\mathbf P(L_p)\times\mathbf P(L_q)$ has centre disjoint from that
Segre variety. A positive-dimensional projective fibre would meet
the centre hyperplane in its ambient fibre space, so fibres are finite.
Normalize a first nonzero coordinate of $C$ to one in finitely many
charts. Gr\"obner bases give zero-dimensional coordinate algebras;
rational roots of coordinate-multiplication characteristic polynomials,
followed by testing the equations, enumerate all rational points.
For each first-factor line choose primitive integral $C_0$. The
solution $D_0$ of $[C_0,D_0]=W$ is unique. Integral factors exist
only if $D_0$ is integral, and are then exactly
\[
 (kC_0,D_0/k),\qquad k\text{ a nonzero signed divisor of the content of }D_0.
\]

For $p=q$, the bracket map $\Lambda^2L_p\to L_{2p}$ is injective:
split associative words into two length-$p$ blocks to recover
$C\otimes D-D\otimes C$. The exterior preimage of $W$ must be
integral and have alternating rank two. Intersect its rational support
with $L_\Z$ to get a saturated rank-two lattice. In an oriented
basis the exterior tensor is $h e_1\wedge e_2$, $h>0$. Enumerate
the finitely many index-$h$ sublattices by Hermite form. Every other
oriented basis is reached by the exact pair moves, leaving higher
coordinates unrestricted. This proves finite completeness of the
leading list, not just the existence of some leading decomposition.

\subsection{Correction kernels and exact periods}
Fix $C,D$, and put $n=c-d$. At offset $s$ the integral correction
map is $A_s(U,V)=[U,D]+[C,V]$ on the full components of weights
$p+s,q+s$. We use the two-coefficient inner-solution theorem:
if $a,b$ belong to a free basis and $[a,V]+[U,b]=0$, then
$U,V\in\Lie(a,b)$~\cite{remeslennikov-stoehr,altassan}. It is
not applied to arbitrary coefficients without first constructing
a free basis containing them.

For $0<s<q-p$, a kernel has dimension at most one. For nonzero
$U$, the graded subalgebra
\[
 \Q C+\Q U+\bigoplus_{j\ge p+s+1}L_j
\]
has a homogeneous free basis containing $C,U$, and contains $D,V$.
The inner-solution theorem puts $D,V$ in $\Lie(C,U)$. Two independent
first components $U_1,U_2$ would similarly give
$D\in\Lie(C,U_1)\cap\Lie(C,U_2)=\Q C$, impossible by weight.
These nonzero one-dimensional kernels are called exceptional.
At $s=q-p>0$ the kernel is $\Q(D,0)$: use
$\Q C+\bigoplus_{j\ge q}L_j$, kill $C$ and higher generators,
and apply the centralizer fact. Its period is the exact substitution
$x\mapsto y^kx$, with all residues retained if $D$ is nonprimitive.
After that offset every kernel pair lies in $\Lie(C,D)$, by applying
the same theorem in
$\Q C+\Q D+\bigoplus_{j\ge q+1}L_j$; for equal weights use
$L_{\ge p}$.

Here is a finite period construction for all these latter kernels.
In the free Lie algebra on $a,b$ with weights $p,q$, truncated at
$c$, set $W_0=\log(e^{-a}e^{-b}e^ae^b)$. There is a filtered
automorphism $\sigma$, with identity associated graded, satisfying
$\sigma([a,b])=W_0$. Successively remove homogeneous discrepancies
by correcting $a,b$: every derived Lie element is in
$[a,H]+[H,b]$, by Jacobi induction. This is the degreewise
symplectic-expansion construction; the related prior construction
is credited in~\cite{kuno} and the retained gauge audit.

If $[u,b]+[a,v]=0$ at offset $s$, the positive-weight derivation
$\delta(a)=u$, $\delta(b)=v$ kills $[a,b]$. Thus
\[
 \alpha=\sigma\exp(\delta)\sigma^{-1}
\]
fixes $W_0$ and has first increments $u,v$. Its powers have rational
polynomial coefficients in the exponent. In the weighted two-generator
Hall group obtained by discarding weights greater than $c$, convert
$\exp(\alpha^k(a)),\exp(\alpha^k(b))$ to group coordinates.
Choose $M>0$ divisible by all denominators of their nonconstant
coefficients. At $k=\pm M$ the coordinates are integral, so
$\alpha^M$ and its inverse give actual group automorphisms fixing
$[A,B]$. Substitution $A\mapsto x$, $B\mapsto y$ gives reversible
pair substitutions in the target group, not necessarily automorphisms
of that target.

Noncommuting $C,D$ freely generate their subalgebra. Express a complete
rational kernel basis in them and apply this construction to each
integer kernel direction. The resulting integral multiples span a
full-rank exact period lattice. Each such layer therefore has finitely
many residue branches. Replacing the entire lattice by one rational
representative would not be valid.

\subsection{A diagonal positional-polynomial lemma}\label{app:positional}
Let $E_0,E_1,\ldots$ freely generate an associative algebra over $\Q$,
with $\delta(E_i)=E_{i+1}$. Encode the coefficient of
$E_{i_1}\cdots E_{i_m}$ as that of $z_1^{i_1}\cdots z_m^{i_m}$.
This is a linear dictionary, not an algebra homomorphism; $\delta$
becomes multiplication by $z_1+\cdots+z_m$.

Suppose $\delta V=[E_0,D]$ and $P$ encodes an $m$-letter $D$.
Then $V$ has encoding
\begin{equation}\label{eq:positional-quotient}
 \frac{P(z_2,\ldots,z_{m+1})-P(z_1,\ldots,z_m)}
 {z_1+\cdots+z_{m+1}}.
\end{equation}
Polynomiality implies cyclic invariance of $P(y_1,\ldots,y_m)$
on the hyperplane $y_0+\cdots+y_m=0$. Let $R$ denote restriction
of an $(m+2)$-letter encoding to
\[
 z_{m+2}=z_1,\qquad z_1+\cdots+z_{m+2}=0.
\]

\begin{lemma}\label{lem:diagonal}
If $R([E_0,V])=0$, then $P$ is constant.
\end{lemma}
\begin{proof}
Put $z=(z_1,\ldots,z_m)$ and $a=-\frac12\sum z_i$.
Substituting~\eqref{eq:positional-quotient} into the assumption gives
\[
 2P(z)=P(z_2,\ldots,z_m,a)+P(a,z_1,\ldots,z_{m-1}).
\]
The calculation may divide by $a$ on the nonempty open set $a\ne0$,
then extends polynomially. Cyclic invariance rewrites it as
\[
 2P(z)=P(z+ae_1)+P(z+ae_m).
\]
On the zero-sum hyperplane, let $T_{i\to j}$ halve coordinate $i$
and add the other half to $j$. The last identity says $P$ is the
average of its values under the two transfers from coordinate zero
to its cyclic neighbors. Cyclic invariance gives the same identity
at each coordinate. All these maps preserve the hyperplane and do
not increase the $\ell^1$ norm.

On a closed $\ell^1$ ball, $P$ attains a maximum. At a maximizing
point both terms of each average must again maximize, so every finite
transfer composition preserves maximality. Repeat a complete clockwise
cycle of transfers $m+1$ times. Its matrix $B$ is positive and
column-stochastic: retained halves permit waiting, and clockwise paths
carry positive mass between every pair of coordinates. If $\epsilon$
is its smallest entry and $\ell=m+1$, write
$B=\epsilon J+(1-\ell\epsilon)C$ with $C$ column-stochastic when
the latter coefficient is positive. On the zero-sum hyperplane
$Jy=0$, so $B$ strictly contracts $\ell^1$; if that coefficient is
zero it annihilates the hyperplane. Iteration tends to zero, making
the maximum $P(0)$. The same applies to the minimum. Every ball
therefore has constant value.
\end{proof}

For a Lie element $D$, a nonzero constant positional polynomial would
be a multiple of $E_0^m$, which is a Lie element only for $m=1$.
Except for $D=\lambda E_0,V=0$, the nonzero Lie solution therefore
has $R([E_0,V])\ne0$. The other two identities needed below are
\begin{equation}\label{eq:r-kills}
 R(\delta F)=0,\qquad R([E_j,F])=z_1^jR([E_0,F]).
\end{equation}
They follow immediately from the restriction and apply at each fixed
letter count. Injectivity of $\delta$ on positive-length tensors is
injectivity of multiplication by their positional sum.

\subsection{Why all later columns fail to cancel the first quadratic}
Fix all pair coordinates below the first exceptional offset $t<q-p$.
If $n\ge2t$, keep all coordinates at offsets $t,\ldots,2t-1$ and
solve their commutator equations jointly. They are an affine integer
linear system: products of two increments first enter at offset $2t$.
Its complete solution set is empty or an affine integer lattice
$\mathcal S$. The first kernel is a line. If its coefficient varies
on $\mathcal S$, choose an integral lattice basis with that coefficient
$k_0+bT$, $b\ne0$, and every other direction starting strictly later.
At offset $2t$, the new error has the form
\[
 q(T)+B\mathbf s=0\pmod{\im A_{2t}},
\]
with $q$ of degree at most two and $B$ constant. The needed statement is
\begin{equation}\label{eq:full-separation}
 q_2\notin\Span_\Q(\im A_{2t},\text{ all columns of }B).
\end{equation}
Showing merely $q_2\notin\im A_{2t}$ would leave an essential gap.

Write $X=\log x$, $Y=\log y$. In the free graded subalgebra
$L_{\ge p}$ include $C$ in a homogeneous free basis. A fixed filtered
substitution, identity on associated graded, removes the fixed
$X$-components at offsets $1,\ldots,t-1$. This is a rational proof
normalization independent of the parameters. Put
\[
 Z=h(\operatorname{ad}_X)Y,\qquad h(z)=(1-e^{-z})/z.
\]
The part of $\log[x,y]$ linear in $Y$ is $[X,Z]$. Every other term
has at least two $Y$ occurrences and one $X$; any dependence on a
block parameter has weight at least $p+2q+t>d+2t$. Such terms may
contribute to the fixed target but do not affect its varying block.
The required components of $X,Y,Z$ are affine in the block parameters:
two varying $X$ increments already have weight $2(p+t)>p+2t$,
and the analogous inequalities handle $Y$ and $h(\operatorname{ad}_X)Y$.

Use the graded free subalgebra
$\mathcal A=\Q C+\bigoplus_{j\ge p+t}L_j$. Include $C$ and the
nonzero first kernel component $U$ in a homogeneous free basis and
project onto $\Lie(C,U)$. The inner-solution theorem puts the leading
$D$ and corresponding $V$ in this subalgebra, so the projection fixes
them. Including all of $L_{p+t}$ also retains arbitrary particular
solutions at that offset; it is not a projection justified only for
the kernel line.

Set $E_j=(\operatorname{ad}_C)^jU$, $\delta=\operatorname{ad}_C$.
Lazard elimination makes the ideal of $U$ free on these $E_j$.
Let $D_m$ be the nonzero component of $D$ with maximal $E$-letter
count. Its associated $V_{m+1}$ satisfies
$\delta V_{m+1}=-[E_0,D_m]$. The harmless sign change in
Lemma~\ref{lem:diagonal} yields
\[
 R([E_0,V_{m+1}])\ne0.
\]
The constant exception is excluded by the Lie condition if $m>1$
and by $q>p+t$ if $m=1$.

Every projected first-factor component at offset $i\le2t$ has at
most one $E$ letter, since $2(p+t)>p+2t$. Thus it is zero or a
multiple of $E_j$ with $pj=i-t$. Consequently
$[L_{p+2t},D]$ has at most $m+1$ such letters, and $R$ kills
$[C,L_{q+2t}]$. At letter count $m+2$ it annihilates $\im A_{2t}$.

For $0\le j<t$, let $F_j$ be the letter-count-$(m+1)$ part of the
fixed $Z$ component at offset $j$; $F_0=0$. The coefficient of $T$
in the already satisfied equation at degree $d+t+j$, followed by
projection and $R$, gives the triangular convolution
\[
 bR([E_0,F_j])+\sum_{i=t+1}^{t+j}
 b_i z_1^{(i-t)/p}R([E_0,F_{t+j-i}])=0.
\]
Only integral exponents occur. Induction gives $R([E_0,F_j])=0$
for all $j<t$. At degree $d+2t$, any other parameter starts later
than $t$: its varying $X$ part pairs only with such fixed $Z$ parts,
and its varying $Z$ part pairs with $C$ or a vanished fixed lower
$X$ part. Equation~\eqref{eq:r-kills} therefore kills every column
of $B$, including universal-kernel columns.

The coefficient of $T^2$ is $b^2[U,V]$, whose letter-count-$(m+2)$
part has nonzero $R$-image. This proves~\eqref{eq:full-separation}.
After the lower equations have been satisfied, the target error begins
in degree $d+2t$. Hall and logarithmic coordinates agree there, and
the filtered normalization is identity in that first error degree.
Thus the functional pulls back to the actual group correction space.
The retained \record{problems/N8/group-block-supplement.md}{coefficient dictionary}
spells out this group interface and the later columns.

\subsection{Termination and completeness}
For each leading pair process successive offsets. A zero kernel has
at most one integral affine solution. A universal kernel has finitely
many branches modulo its full-rank exact period lattice. At the first
exceptional offset $t$, if fewer than $2t$ offsets remain from the
leading degree, the whole remaining system is jointly linear and is
decided at once.

Otherwise form $\mathcal S$. Reject an empty lattice. If its first
exceptional coefficient is constant, fix the prefix through $t$ and
restart at $t+1$. If it varies, use rational linear algebra to find a
functional annihilating $\im A_{2t}$ and all later columns but not
$q_2$, as guaranteed by~\eqref{eq:full-separation}. Its scalar
quadratic has at most two integral roots. For each root fix the entire
prefix through $t$ and restart, freeing higher coordinates again.

Freeing those coordinates may repeat unsuccessful continuations, but
cannot lose a solution. A positive answer still requires every equation
through class $c$ and a final actual commutator check. The nonunit step
$b$, all integer congruences and all exact-period residues remain in
the corresponding integral systems. A rational correction or rational
root is never accepted as an integral one. Each restart increases the
fixed offset and every branch has finitely many children, so the tree
terminates. Conversely any solution has a normalized leading branch
and survives every necessary equation, exact period and quadratic root
choice. This proves the proposed decision and construction.

The proof needs no bound on the number of exceptional offsets and no
integer-curve oracle. The positional argument has substantial Lean
support, including its ordered-word dictionary and full correction-span
formulation. The abstract free-Lie identification, the production of
the group block and the full recursive program are outside Lean.
Higher-weight positive and negative native group controls support
specific instances; they do not establish the structural proof in
all ranks and classes.
